你按论文里的公式手写了一层的反向，跑通了，损失也在降。但降得比别人慢一截。

问题最后出在一个转置上。公式没抄错，是布局约定不同——论文用分母布局，你的框架用分子布局，两者的结果互为转置。方阵掩盖了这个错误，梯度检查也没报警，因为形状对得上。

这类事故在矩阵求导里占了大多数，而它们几乎都不是微积分问题。真正的困难有两个：一是\emph{形状与约定}，二是\emph{怎样在不逐元素展开的前提下把梯度整理出来}。这一部分解决的就是这两件事。

底下的想法只有一句：\textbf{导数首先是一个线性映射，矩阵只是它在选定基下的记账方式。}把这句话当真，很多东西会自动归位——梯度是行是列取决于你怎么记账，链式法则是映射复合而矩阵乘法只是它的坐标表示，而\emph{乘法的结合顺序决定计算代价}这件事，就是反向传播为什么高效的完整回答。

\begin{center}
\begin{tikzpicture}[x=1cm,y=1cm,font=\small,>=stealth]
  \draw[black!45,fill=blue!9] (0,1.35) rectangle (3.6,2.85);
  \node[align=center] at (1.8,2.1) {导数是\\线性映射};
  \draw[black!45,fill=green!8] (5.1,3.45) rectangle (13.3,4.35);
  \node[align=left,anchor=west] at (5.3,3.9) {形状与布局：先对形状，再对公式（第一章）};
  \draw[black!45,fill=green!8] (5.1,2.3) rectangle (13.3,3.2);
  \node[align=left,anchor=west] at (5.3,2.75) {微分法：整理成 \(\ip{G}{\mathrm dX}\)，读出 \(G\)（第二章）};
  \draw[black!45,fill=green!8] (5.1,1.15) rectangle (13.3,2.05);
  \node[align=left,anchor=west] at (5.3,1.6) {trace 演算：让矩阵表达式可以重排（第三章）};
  \draw[black!45,fill=green!8] (5.1,0.0) rectangle (13.3,0.9);
  \node[align=left,anchor=west] at (5.3,0.45) {落到计算图：从右往左乘就是反向传播（第四章）};
  \foreach \y in {3.9,2.75,1.6,0.45}
    \draw[->,black!55] (3.65,2.1) -- (5.05,\y);
\end{tikzpicture}

\emph{四章是同一个想法的四层落地。最后一层解释框架每天在做什么，前三层解释你为什么能在它算错时看出来。}
\end{center}

\subsection{四条贯穿始终的判断}

\textbf{先对形状，再对公式。}输入 \(m\times n\)、输出标量，梯度必然是 \(m\times n\)。这个检查不需要任何推导，却能拦下大部分错误。方阵是它的盲区——梯度检查时特意让各维长度互不相同，转置错误立刻暴露。

\textbf{不要逐元素展开。}写出 \(\mathrm dL=\ip{G}{\mathrm dX}\) 的形式，\(G\) 就是梯度。求导于是变成一件代数整理工作：用 trace 的循环性把表达式重排，直到微分项被孤立出来。推一个新层的梯度，走这条路比写下标快得多，也不容易错。

\textbf{数学公式与数值实现是两笔账。}\(\mathrm d(X^{-1})=-X^{-1}(\mathrm dX)X^{-1}\) 是正确的公式，但实现时绝不该显式构造逆矩阵，应写成解线性系统；\(\mathrm d\log\det X=\trace(X^{-1}\mathrm dX)\) 同样正确，而 \(\log\det\) 要从 Cholesky 的对角线取对数求和，先算行列式会直接溢出。这两条在数值分析那一部分有更完整的说明。

\textbf{乘法顺序就是模式选择。}链式法则给出一串矩阵相乘。从右往左（从输出端起步）每步都是行向量乘矩阵，代价 \(O(d^2)\)；从左往右则是矩阵乘矩阵，代价 \(O(d^3)\)。反向传播就是前一种顺序，它的代价与前向同阶——这不是工程技巧，是结合律的直接后果。反过来，输入维数远小于输出维数时，前向模式才是划算的那个。

\subsection{四章怎么安排}

\hyperref[ch:081]{``导数的对象与记号''}篇幅最短却最该先读，它建立形状检查这个动作。\hyperref[ch:082]{``微分、链式法则与二阶信息''}是主力工具，其中链式法则那一篇解释了反向传播的代价结构，与深度学习部分直接衔接。

\hyperref[ch:083]{``Trace 方法与常用矩阵导数''}更像手册，可以按需查阅；但其中逆矩阵与 \(\log\det\) 两篇请读完，它们讲的是公式与实现的分离，实践价值高于推导本身。

\hyperref[ch:084]{``从梯度推导到反向传播''}把前三章落地：最小二乘与 logistic 的梯度用同一套方法推出，softmax 与交叉熵的化简解释了框架为何把两者合成一个算子，最后一篇把整套东西放到计算图上。

手推一遍的意义不在于以后都手推——框架会替你做。意义在于，当曲线降得不对劲时，你知道该去看哪一步。
